% !TeX program = xelatex
% !BIB program = biber
% 通用中文学术综述版式；非期刊官方投稿模板。
\documentclass[UTF8,a4paper,zihao=-4,fontset=fandol]{ctexart}
\usepackage[left=22mm,right=22mm,top=22mm,bottom=22mm,headheight=16pt,headsep=7mm,footskip=10mm]{geometry}
\usepackage{graphicx,xcolor,caption,fancyhdr,needspace,placeins}
\usepackage{amsmath,amssymb,booktabs,tabularx,array,longtable}
\usepackage[backend=biber,style=gb7714-2015,sorting=none,gbnamefmt=uppercase,gbpub=false,doi=true,url=true,isbn=false]{biblatex}
\addbibresource{references.bib}
\AtBeginBibliography{\interlinepenalty=10000}
\usepackage[unicode,hidelinks,pdfauthor={范举；范梅浩},pdftitle={智能数据准备研究综述：关键技术、流程优化与评测方法}]{hyperref}
\definecolor{ink}{HTML}{173753}
\definecolor{muted}{HTML}{647580}
\definecolor{rulegray}{HTML}{CCD7DF}
\ctexset{
 section={format=\large\bfseries,number=\arabic{section},aftername=\hspace{0.6em},beforeskip=16pt,afterskip=8pt},
 subsection={format=\normalsize\bfseries,number=\thesection.\arabic{subsection},beforeskip=11pt,afterskip=6pt}
}
\AtBeginDocument{\linespread{1.0}\fontsize{11pt}{17pt}\selectfont\special{pdf:minorversion 7}}
\setlength{\parindent}{2em}
\setlength{\parskip}{2pt}
\setlength{\emergencystretch}{2em}
\setlength{\textfloatsep}{13pt plus 3pt minus 3pt}
\setlength{\intextsep}{12pt plus 3pt minus 3pt}
\renewcommand{\topfraction}{0.92}
\renewcommand{\bottomfraction}{0.86}
\renewcommand{\textfraction}{0.08}
\renewcommand{\floatpagefraction}{0.75}
\setcounter{topnumber}{2}
\setcounter{bottomnumber}{2}
\setcounter{totalnumber}{4}
\captionsetup{font=small,labelfont=bf,labelsep=quad,justification=justified,singlelinecheck=false,skip=6pt}
\pagestyle{fancy}
\fancyhf{}
\fancyhead[L]{\small 智能数据准备研究综述}
\fancyhead[R]{\small 范举，范梅浩}
\fancyfoot[C]{\small\thepage}
\renewcommand{\headrulewidth}{0.4pt}
\newcolumntype{Y}{>{\raggedright\arraybackslash}X}
\newcolumntype{L}[1]{>{\raggedright\arraybackslash}p{#1}}
\renewcommand{\arraystretch}{1.12}
\newcommand{\articlefigure}[3]{%
\begin{figure}[htbp]
\centering
\includegraphics[width=\linewidth]{figures/#1.png}
\caption{#2}\label{#3}
\end{figure}}
\begin{document}
\thispagestyle{plain}
\begin{center}
{\fontsize{20pt}{27pt}\selectfont\bfseries 智能数据准备研究综述\par}
\vspace{3pt}
{\fontsize{15pt}{21pt}\selectfont\bfseries 关键技术、流程优化与评测方法\par}
\vspace{11pt}
{\large 范举\quad 范梅浩\par}
\vspace{5pt}
{\small 中国人民大学信息学院，数据工程与知识工程教育部重点实验室\par
北京 100872\par}
\end{center}
\vspace{5pt}
% 正文按完整句子换行。参考文献、图注、作者简介不计入正文篇幅。
一家连锁商店想比较各门店的净销售额，看起来只需做一次加减法。
但订单、退款和门店信息往往分散在不同表格中，日期写法不一，同一家门店可能使用不同简称，一笔订单还可能对应多笔退款。
假设一笔100元的订单分两次退款，金额分别为20元和10元。
如果把订单与退款记录直接关联，再对每行相减求和，就可能把100元算两遍，得到170元，而不是正确的70元。
程序没有报错，图表也很漂亮，结论却已经偏离事实。

分析之前，必须先弄清数据的含义、组织方式和使用规则。
找到相关数据，识别问题，连接不同来源，完成必要转换，再检查结果，这些工作共同构成了数据准备。
它不是分析过程中的附属杂务，而是把原始记录变成可靠依据的关键环节。
智能数据准备希望进一步降低这项工作的门槛：让人说明目标，让系统协助理解数据、设计流程、执行检查，并在有疑问时寻求确认，而不是要求使用者从头编写每一行程序。\cite{demystify}

\section{数据准备，准备的究竟是什么}

数据清洗只是数据准备的一部分。
清洗主要处理重复、缺失、错误和冲突；准备还包括发现数据、解释字段、整合来源、改变组织形式，以及为学习任务补充必要的标注。
例如，将多个月份分别占一列的报表改成按月份逐行记录，是结构转换；判断两个不同名称是否指向同一家企业，是实体匹配；把订单与退款组织到同一个统计口径下，则需要同时考虑关联与计算。
这些操作往往相互影响，不能机械地按固定顺序执行。\cite{demystify,quality}

准备对象也不局限于电子表格。
合同里的条款、图片中的文字、设备日志中的时间记录，常需要先提取并组织，才能与数据库中的信息共同使用。
提取过程本身也可能出错，因此既要保存整理后的结果，也要保留能够返回原文、原图或原始记录的位置。
这使数据准备同时承担了连接不同信息形式和保留证据的任务。\cite{demystify,datasheets}

准备也不是分析前只做一次的工作。
当源系统增加字段、门店改变名称，或统计规则发生调整，原有流程可能不再适用。
可持续的数据准备需要识别这些变化，重新检查受影响的步骤，并保留不同版本之间的关系。
否则，一次正确的结果并不能保证下个月仍然正确。\cite{quality}

“准备好”也没有脱离用途的统一答案。
分析门店经营，需要保留退货信息；识别异常交易，恰恰不能先删掉看起来不正常的订单。
同一个空白值，可能表示尚未采集、不适用，也可能意味着系统漏记，不能一律填成零。
数据质量因此既包括准确性、完整性和一致性，也取决于它是否及时、是否适合当前任务。
把表格整理得整齐，并不等于把事实解释正确。\cite{quality}

智能之处首先在于理解这种“用途与处理方式”的关系。
传统程序按明确规则执行，语言模型则能够利用列名、文本内容和任务描述判断哪些操作可能需要进行。
但语言模型的判断并不是事实本身。
比较可靠的分工是：模型负责理解与提出方案，数据库和经过检验的程序负责运算，质量检查负责发现偏差，人负责确认业务含义和重要决策。
因此，给工具加上聊天窗口只是入口，能否根据真实执行结果修正方案，才是更重要的能力。\cite{fm,autoprep,deepprep}

\articlefigure{fig01_scope}{从原始数据到可信结果。数据准备围绕具体用途展开，发现、理解、清洗、整合、转换与验证可以反复交织，来源和变更记录贯穿全过程。}{fig:scope}

\section{它与我们的生活和工作有什么关系}

\textbf{在办公与经营中，减少重复劳动只是第一步。}
合并月报、核对库存、分析客户反馈，都需要处理来源不同、结构不齐的记录。
智能系统可以建议日期转换、字段对应和统计步骤，把一次处理保存为可重用的流程。
真正重要的收益，是让使用者有机会看清计算依据，而不是仅仅更快得到一个数字。
开篇的退款例子中，系统应先确认统计规则，再判断是否需要按订单汇总退款，最后检查销售额是否被重复计算。
这也是自然语言驱动的数据准备研究试图解决的实际问题。\cite{autoprep,deepprep}

\textbf{在科研和公共服务中，准备数据要保留解释所需的背景。}
把不同设备的观测记录放在一起，不能只统一文件格式，还要确认单位、采样时间和测量条件。
罕见的高值可能是传感器故障，也可能是真实的极端事件，删除之前需要查证。
医疗研究中的不同编码和记录方式同样需要谨慎对应，自动处理不能替代专业判断。
科学数据管理强调可查找、可访问、可互操作和可重用，背后的要求正是让后来使用数据的人理解它的来源、含义与适用范围。\cite{fair}

\textbf{在大模型训练中，数据准备影响模型学到什么。}
网页、文档、图片和问答材料需要去重、筛选、整理来源，并检查不同内容的覆盖情况。
只追求数量，可能保留大量重复信息；过度追求“干净”，又可能删去少见但有价值的表达。
Data-Juicer等系统把处理步骤组织成可组合的操作，并连接数据分析和模型评估，使研究者能够比较不同处理方案，而不只是凭经验选择材料。\cite{datajuicer,llmdata}
这里有两个不同方向：用人工智能帮助准备数据，以及为人工智能准备数据。
两者相互促进，但不能混为一谈。

\textbf{在机器人学习中，还要保证不同记录描述的是同一时刻、同一个动作。}
相机画面、机械臂位置和控制指令即使各自正确，时间没有对齐，也可能形成错误的学习样本。
不同机器人还可能使用不同的动作表示。
Open X-Embodiment对多种机器人数据进行汇集和标准化，说明跨设备学习离不开数据组织工作。\cite{oxe}
让智能系统进一步协助检查时间同步、动作含义和样本覆盖，是值得探索的方向，而不是已经普遍解决的问题。
图\ref{fig:applications}概括了这些用途，它们的共同点是：数据准备必须服务于具体任务，而非套用同一套清洗规则。

\articlefigure{fig02_applications}{数据准备服务于不同用途。办公经营关注口径和关联，科研关注来源与测量条件，大模型训练关注质量与覆盖，机器人学习还需要时间和动作的一致性。图中为用途示意，并非特定系统的部署案例。}{fig:applications}

\section{研究正在从单项工具走向完整流程}

\subsection{让机器帮助人整理数据}

数据准备并不是大模型出现后才开始的研究。
早期系统主要依靠规则与脚本，随后出现了更易操作的交互式工具。
Wrangler让用户通过表格界面示范或选择变换，再生成相应处理步骤，降低了逐行编程的负担。\cite{wrangler}
HoloClean进一步综合数据约束、统计关系和参考信息推断可能的修复结果，说明机器能够帮助发现仅靠手写规则难以覆盖的问题。\cite{holoclean}

另一条路线是自动寻找转换程序。
Auto-Tables针对结构不规则的表格合成多步转换，并不要求用户先给出转换后的示例表。\cite{autotables}
这类研究提醒我们，自动化不只有“大模型写代码”一种实现方式。
规则、程序搜索和统计学习各有所长，至今仍是智能系统的重要组成部分。

语言模型带来的新变化，是能够较自然地理解字段含义和任务描述。
早期探索已经将实体匹配、错误检测等数据任务转化为语言问答；Jellyfish则通过针对性训练，使可本地运行的模型支持多类数据预处理任务。\cite{fm,jellyfish}
不过，做好一次匹配或一次修复，不等于能稳定完成一条相互依赖的长流程。
研究的重心因而开始从单项能力转向多步骤协同。

\subsection{从理解问题，到根据执行结果回退修正}

AutoPrep关注一个容易被忽视的事实：同一张表面对不同问题，需要的准备工作并不相同。
它把任务分为规划、编程和执行，让系统先根据问题确定需要哪些高层操作，再转成具体程序。
例如，回答某年的年龄分布可能需要从出生日期计算年龄，而回答地区分布则需要识别和规范化地域信息。
这项工作的价值，是把数据准备与实际问题连接起来，而非不加区分地清洗所有字段。\cite{autoprep}

DeepPrep进一步处理多表、多步骤准备中“错误很晚才暴露”的困难。
它保留中间表和操作反馈，通过树状的过程组织方式保存不同尝试，使系统可以退回较早的决策，而不只是反复修改最后一段代码。
它还通过分阶段训练，让模型学习操作使用、执行反馈和回退行为。
在论文所测任务中，这种专门训练改善了开源模型的准确率与推理成本权衡，但这不意味着任意业务数据都能无人值守地处理。\cite{deepprep}

图\ref{fig:agent}用退款例子说明这种区别。
当系统发现销售额重复，真正需要修正的可能是前面的关联方式，而不是后面的加减公式。
能够保存现场、找到较早的错误步骤并重新尝试，是数据准备智能体区别于一次性代码生成的重要方向。
这里的“智能体”，指能够围绕目标调用工具、观察结果并继续行动的系统，不意味着它已经拥有人的全部判断能力。

\articlefigure{fig03_agent}{执行反馈与回退修正。示例中，一笔100元订单对应20元和10元两笔退款。直接关联后求和可能得到错误的170元；先汇总退款、再关联计算，才能得到70元。数字为解释原理而设，不是实验结果。}{fig:agent}

\subsection{让方法适应新数据，也让能力能够被检验}

另一项挑战是换一个行业、换一批表格后，系统能否继续工作。
KnowTrans通过复用已有任务的知识，并引入外部信息，帮助专用数据准备模型以较少样本适应新数据和新任务。
它并非证明模型已经无需领域知识，恰恰说明恰当的知识补充仍然重要。\cite{knowtrans}

Text-to-Pipeline把自然语言要求转化为处理流程作为明确任务，并构建PARROT基准，包含约1.8万项任务和16类核心操作。
其研究发现，模型不仅会弄错操作之间的组合关系，也会把文字要求对应到错误的列、条件或参数。
配套的Pipeline-Agent利用中间执行状态改进流程，但仍存在明显错误。\cite{texttopipeline}
BAT则探索在没有目标表实例的条件下，通过受控操作和树搜索合成准备流程。\cite{bat}
这些工作与前述研究从不同入口推进同一问题：如何把大致合理的意图，落实成可以检查的操作。

总体而言，现阶段既不是“数据准备已经全面自动化”，也不是“模型只能做演示”。
单项任务已有丰富方法，较明确的多步骤任务正在取得进展，跨领域、含糊需求和长期运行仍是薄弱环节。
PrepBench将意图澄清、代码生成和工作流表达纳入评测，也发现现有系统仍难稳定贯通这些能力。\cite{prepbench}
不同论文采用的数据、任务和成本口径并不相同，不能把各自的分数直接排成一张通用排行榜。

\section{从能运行到可信赖，还需要跨过哪些门槛}

\textbf{首先是说清目标，而不是让模型猜口径。}
“统计活跃用户”究竟按登录、购买还是互动计算？
“删除重复”是删除完全相同的行，还是合并同一对象的多条记录？
这些歧义会改变结果。
在动手转换之前，还要确认选中的数据是否来自正确的时间段，字段说明是否对应当前版本。
数据目录、业务词典和经过确认的历史流程可以提供帮助，但系统应展示采用了哪份说明，避免把过时经验直接套到新任务上。
系统应该优先追问少数影响最大的条件，并在执行前展示拟采用的定义。
PrepBench的结果表明，交互澄清能够帮助改进表现，但问题问得多不等于问得好。\cite{prepbench}

\textbf{其次是把纠错建立在可观察的证据上。}
程序成功运行，只能说明它完成了计算，不能说明计算符合原意。
关联前后记录数为何变化，关键标识是否仍然唯一，某类样本是否被大量过滤，都应进入检查过程。
中间状态和变更记录可以帮助追溯问题，而独立的规则、统计检查与人工抽查应与模型判断相互补充。
执行反馈和回退机制为此提供了基础，但验证标准本身仍需要领域知识。\cite{quality,deepprep}

\textbf{第三是承认无法恢复的信息。}
模型可以根据上下文提出缺失值的候选答案，却不能保证猜测等于真实记录。
对难以确认的内容，应保留缺失标记或多个候选，并说明依据，而不是把不确定性藏进一张看似完整的表里。
处理异常值也应区分错误与少见事件，避免“修好”数据的同时抹掉研究对象。
数据集说明文档所强调的来源、组成、处理过程和使用限制，有助于把这些边界传递给后续使用者。\cite{datasheets}

\textbf{第四是让可靠性与规模、成本一起设计。}
面对海量记录，不宜让语言模型逐行处理所有内容。
可以先用数据库完成计数、分组和确定性的格式转换，再把难以解释的字段、少量代表样本和检查摘要交给模型。
已确认的步骤可以缓存和重用，数据变化后只重新处理受影响的部分。
但抽样可能漏掉稀有错误，因此必须配合覆盖全量数据的规则检查和有针对性的复核。
这种分工比单纯扩大模型更符合数据系统的特点。\cite{demystify,datajuicer}

\textbf{第五是把权限和责任放进流程。}
系统不应因为“能够调用工具”，就默认可以读取全部记录、向外部服务发送敏感内容或覆盖原始数据。
本地运行可以减少数据外传的路径，却不能替代访问控制和操作审计。\cite{jellyfish}
合理的设计是保留只读原始副本，限制执行权限，记录关键修改，并把不可逆操作交给人确认。
交付物也不应只有结果表，还应包括处理流程、验证依据和适用范围。

\textbf{最后是评价真正的使用效果。}
好的准备流程不仅要生成正确结果，还要控制时间、计算开销和人工修订量。
面向模型训练，还应检查不同数据群体的覆盖和下游表现；面向经营分析，则要检查口径变化是否影响结论。
评测时必须避免把测试答案或未来信息混入准备过程，也不能用总体分数掩盖少数高风险失败。
基准测试有助于比较方法，真实应用仍需要结合自己的数据、风险和长期维护成本做验证。\cite{prepbench,llmdata}
图\ref{fig:trust}将这些要求归纳为相互配合的检查环节。

\articlefigure{fig04_trust}{可信数据准备的关键检查。目标、执行、事实、成本、权限与效果需要共同约束系统行为；自动化程度应随风险调整，而不是把所有任务一律交给模型。}{fig:trust}

\section{未来：从一次性整理走向持续的数据协作}

更可行的发展路径，是逐步扩大经过验证的自动化范围。
低风险、规则明确、重复发生的工作，可以更多交给系统执行；涉及口径变化、敏感信息和重要决策的步骤，则保留人的确认。
这并非自动化的退让，而是把人的专业判断用在最值得介入的地方。
能够说明何时需要帮助，与能够独立完成任务同样重要。

另一个方向，是让准备过程与后续分析相互反馈。
当分析发现某类记录持续缺失，或模型在某些场景中频繁失败，系统可以重新检查数据覆盖、标注方式和处理规则。
DeepAnalyze等数据科学智能体已尝试连接数据处理、分析与报告生成，为这种衔接提供了探索。\cite{deepanalyze}
但更长的流程也意味着错误可能传播得更远，因此每个阶段仍需要独立检查，不能让一份表达流畅的报告替代对数据的核实。

随着多模态和跨设备应用发展，数据准备还需要更好地描述时间、空间、动作和来源之间的关系，并在数据持续更新时保持一致。
可复用的操作定义、数据说明和公开评测，有望减少重复建设，使不同系统的结果更容易交换和核验。\cite{fair,oxe,datasheets}
这些是需要继续推进的研究方向，而不是对短期全面无人化的承诺。

智能数据准备最终要解决的，不只是“让机器替人整理数据”，而是让数据的用途更明确、处理更可靠、错误更容易被发现。
当系统能够同时交付可用的数据、可重现的过程和可检查的依据，人工智能才获得了更值得信任的数据基础。

\FloatBarrier
\clearpage
{\small
\setlength{\bibitemsep}{4pt}
\renewcommand{\bibfont}{\fontsize{9.5pt}{13.2pt}\selectfont}
\printbibliography[title={参考文献}]
}
% 作者简介素材保留于 author_bios.tex，按目标期刊要求单独提供。
% % 作者简介供本人确认；本次未取得可确认授权的两位作者照片。
\par\needspace{62mm}
\section*{作者简介}
\noindent
\begin{minipage}[t]{0.20\linewidth}
\vspace{0pt}\centering\authorphoto{fan_ju}{范举}
\end{minipage}\hfill
\begin{minipage}[t]{0.76\linewidth}
\vspace{0pt}
\textbf{范举}，中国人民大学信息学院教授、博士生导师，国家级青年人才，RUC-DataLab团队负责人。
主要研究数据治理技术与系统、智能数据库系统及数据与人工智能的交叉问题，关注数据准备、数据质量管理和数据科学智能体。
主持国家自然科学基金优秀青年基金、重点项目等，致力于将语义理解、自动化处理与可靠的数据系统结合，提升数据的可用性与应用价值。
\end{minipage}
\par\vspace{17pt}\needspace{48mm}
\noindent
\begin{minipage}[t]{0.20\linewidth}
\vspace{0pt}\centering\authorphoto{fan_meihao}{范梅浩}
\end{minipage}\hfill
\begin{minipage}[t]{0.76\linewidth}
\vspace{0pt}
\textbf{范梅浩}，中国人民大学信息学院2023级直博生，师从范举教授。
主要研究面向数据分析的大语言模型与智能体技术，重点关注任务感知的数据准备、表格理解及多步骤流程的自动生成与验证。
作为第一作者完成AutoPrep和DeepPrep等工作，相关成果发表于VLDB、ICDE等国际会议。
当前关注多模态数据准备以及数据处理过程与下游人工智能任务之间的协同优化。
\end{minipage}

\end{document}
